\begin{tabular}{@{}p{2.5cm}p{4.4cm}p{4.5cm}p{5.4cm}@{}}
\toprule
Finding & Our numbers & What it means & Compared with real networks or theory \\
\midrule
\textbf{A small world: short paths and strong clustering} & average path 6.1 steps (random networks with the same degrees 4.5); average clustering 0.21 vs 0.011 (20$\times$ higher); transitivity 0.21 & Faculty sit in tight groups, yet a few links between groups keep everyone within a few steps. & Typical of co-authorship networks: in seven large ones the average distance is 4.0--9.7 and the clustering (transitivity) 0.072--0.726 (Newman 2001). Ours (6.1 and 0.21) falls inside both ranges. \\
\addlinespace[1.5pt]
\textbf{Well-connected faculty work with each other (assortative)} & degree assortativity r = 0.20 vs $-$0.02 in rewired networks (z = 6.9); by branch r = 0.66, by NIT r = 0.94, by rank only 0.07 & Hubs link to hubs, people stay in their branch and NIT, and rank matters little. & This is the usual pattern of social networks: co-authorship networks have r = 0.12--0.36, while technological and biological networks are disassortative (Newman 2002). \\
\addlinespace[1.5pt]
\textbf{No clean power law in the degree distributions} & a log-normal fits better than a power law for 4 of 4 degree and citation distributions (heavy tails, $\alpha$ $\approx$ 2.5--2.9) & A few people have many co-authors, but the network is not 'scale-free'. & Newman (2001) also found that collaboration networks are not pure power laws: their degree distributions have a cut-off. \\
\addlinespace[1.5pt]
\textbf{The network densifies: links grow faster than people} & links $\propto$ (connected faculty)$^{a}$ with a = 1.65 (R$^2$ = 0.998); the same exponent at the six South Indian NITs is 1.19--1.74 & Each faculty member gains co-authors over time, so the average number of co-authors keeps rising. & Real growing networks densify the same way, with a = 1.08--1.68 (citation graphs 1.56 and 1.68, US patents 1.66, internet routers 1.18, author--paper graphs 1.08--1.15; Leskovec et al. 2005). Ours and five of the six NITs are inside this range; Surathkal is just above it. \\
\addlinespace[1.5pt]
\textbf{Distances shrink as the network grows} & effective diameter 19 steps in 2018 $\to$ 9; average path in the largest group 11.2 $\to$ 6.1; co-authors per connected faculty member 1.67 $\to$ 3.48, but within any one year only 1.7--2.0 & The giant group formed in 2018 through thin bridges (long paths); new links then acted as shortcuts. People do not write with more colleagues each year; the totals rise because old links stay. & Leskovec et al. (2005) found shrinking diameters in real growing networks, largest when the giant component first appears; Barab\'asi et al. (2002) found that in co-authorship networks the separation falls and the average degree rises over time. \\
\addlinespace[1.5pt]
\textbf{Rich get richer: weak for collaborators, stronger for citations} & new faculty co-authors $\propto$ (current co-authors)$^{0.35}$ (Trichy 0.26, Surathkal 0.32); new citations $\propto$ (current citations)$^{0.75}$ & Well-connected faculty gain new collaborators only a little faster than others; highly cited papers gain citations much faster. & Barab\'asi et al. (2002) measured a stronger attachment for co-authorship, k$^{0.8}$ (mathematics) and k$^{0.75}$ (neuroscience). Our collaborator effect is weaker; our citation effect is close to theirs. \\
\addlinespace[1.5pt]
\textbf{New collaborations can be predicted, but exact hits are rare} & a model trained once on 2016--2022 and given only the 2022 network: AUC 0.90 / 0.87 / 0.90 for 2023, 2024, 2025; its top 1\% of pairs (about 1,569) holds 19 / 14 / 9 of the real new pairs, 15--22$\times$ better than random (42 correct matches listed) & Who will work together next is fairly predictable (same department, active people, common colleagues), but only about 1 candidate pair in 2,000 starts each year, so most predicted pairs do not happen yet. & Real co-authorship networks behave the same: a random prediction is right 0.15--0.48\% of the time and the best methods do up to about 55$\times$ better than random, with low raw precision (Liben-Nowell \& Kleinberg 2007). \\
\addlinespace[1.5pt]
\textbf{Ideas spread more easily at Surathkal, as the theory predicts} & SIR epidemic threshold $\beta_c$ = $1/\lambda_{NB}$ = 0.211 (Trichy) and 0.156 (Surathkal); simulated reach from 5 random starters at $\beta$ = 0.1 / 0.2 / 0.3: Trichy 1.9 / 3.9 / 10.1\%, Surathkal 2.8 / 8.4 / 24.8\% & Spreading takes off between $\beta$ = 0.2 and 0.3 at Trichy and around 0.15--0.2 at Surathkal, where the theory puts the thresholds; Surathkal's bigger, connected hubs raise $\lambda$ and so lower its threshold. & The threshold formula for SIR on real networks is the non-backtracking one (Karrer, Newman \& Zdeborov\'a 2014); worked out step by step in our PDF. \\
\addlinespace[1.5pt]
\textbf{Synchrony starts where theory says, but full synchrony is hard} & onset $K_c$ = $K_0/\lambda_1$ = 0.264 (Trichy) and 0.203 (Surathkal); simulated order parameter r at K = 0.2 / 0.3 / 0.4: Trichy 0.08 / 0.11 / 0.16, Surathkal 0.11 / 0.15 / 0.24 (noise level about 0.06); Laplacian eigenratio 670 vs 130 in rewired networks & Partial synchrony begins as predicted, but groups are joined by thin bridges, so the whole network is about 5$\times$ harder to bring into one rhythm than a random one. & The onset formula is Restrepo, Ott \& Hunt (2005); details in our PDF. \\
\addlinespace[1.5pt]
\textbf{Most citations between our own papers are self-citations, at every NIT} & 86\% of the 23,822 citations between our papers share an author; without them none of the old top-10 PageRank papers stays in the top 10; the share is 85--92\% at all six South Indian NITs & The internal citation network mainly shows groups building on their own work, so PageRank inside it is not a measure of influence. & Compared across our six NITs only (no published benchmark). Citing your own earlier work is normal when it is relevant. \\
\addlinespace[1.5pt]
\textbf{Every South Indian NIT densifies; recent hiring goes with a lower exponent, but the evidence is thin} & a = 1.19--1.74 at the six NITs; share hired in 2021--25 against a: $\rho$ = $-$1.00 (exact p = 0.003, n = 6); with a second, independently cleaned dataset for Warangal and Calicut $\rho$ = $-$0.60 (p = 0.24) & Every NIT's faculty network densifies like a real growing network. NITs that hired more recently tend to have lower exponents, but with six NITs this is a pattern, not a law. & All exponents lie near the 1.08--1.68 range of real growing networks (Leskovec et al. 2005). \\
\bottomrule
\end{tabular}
